
\documentclass[11pt]{article}
\usepackage[utf8]{inputenc}
\usepackage{amsmath,amssymb}
\usepackage{enumitem}
\usepackage{geometry}
\geometry{margin=1in}

\begin{document}

\begin{center}
	{\LARGE New York University}\\[6pt]
	{\large Computer Science and Engineering}\\[12pt]
	{\Large CS-GY 6763: Homework 1}\\[6pt]
	{\normalsize Due Wednesday, Feb. 4th, 2026, 11:59pm ET}
\end{center}

\noindent Collaboration is allowed on this problem set, but solutions must be written-up individually. Please list collaborators for each problem separately, or write ``No Collaborators'' if you worked alone.

\section*{Problem 1: Short(er) problems. (10 pts)}
\begin{enumerate}[label=\arabic*.]
	\item (2 pts) Prove Chebyshev's inequality using Markov's Inequality.
\[
\Pr(X \ge t) \le \frac{\mathbb{E}[X]}{t}
\qquad\text{to get }\Pr\!\big[X-\mathbb{E}[X]\ge k\sigma\big]\le \frac{1}{k^2}.
\]

Let
\[
S = (X-\mathbb{E}[X])^2
\qquad \cdots \qquad
\mathrm{Var}(X)=S
\]

Then
\[
\Pr(S \ge t) \le \frac{\mathbb{E}[S]}{t}
\]

\[
\Pr\!\big((X-\mathbb{E}[X])^2 \ge t\big) \le \frac{\mathrm{Var}[X]}{t}
\]

\[
t = k^2\sigma^2
\]

\[
\Pr\!\big((X-\mathbb{E}[X])^2 \ge k^2\sigma^2\big)
\le
\frac{\mathrm{Var}[X]}{k^2\sigma^2}
\]

\[
\text{-- square root the LHS --}
\]

\[
\Pr\!\big(\,|X-\mathbb{E}[X]| \ge k\sigma\,\big)
\le
\frac{\mathrm{Var}[X]}{k^2\sigma^2},
\qquad
\mathrm{Var}(X)=\sigma^2
\]

\[
\Pr\!\big(\,|X-\mathbb{E}[X]| \ge k\sigma\,\big)
\le
\frac{1}{k^2}
\qquad \Box
\]

	\item (4 pts) Consider inserting $m$ keys into a hash table of size $n=5m^2$ using a uniformly random hash function. By the mark-and-recapture analysis from class, we know that the expected number of collisions in the table is
	\[ \frac{m\cdot(m-1)}{2\cdot 5 m^2} < \tfrac{1}{10}. \]
	So, by Markov's inequality, with probability $>9/10$, the table has no collisions (<1 collisions). Thus, we can look up items from the table in worst-case $O(1)$ time. Give an alternative proof of the fact that we have no collisions with $>9/10$ probability in a table of size $c m^2$ for some sufficiently large constant $c$. Specifically, to have no collisions, we must have the following events all happen in sequence: the second item inserted into the hash table doesn't collide with an existing item, the third item inserted doesn't collide with an existing item, \dots, the $m$th item inserted doesn't collide with an existing item. Analyze the probability these events all happen. Hint:\ You might want to use the fact that $\left(1-\tfrac{1}{n}\right)^n \le e^{-1}\le \left(1-\tfrac{1}{n}\right)^{n-1}$ for any integer $n\ge2$.

	\[
	\Pr(\text{No collision inserting $m$ keys})
	=
	1-\Pr(\text{collision inserting $m$ keys})
	\]

	\[
	n = c m^2,\qquad c\in\mathbb{R}
	\]

	\[
	\Pr(h(x)=h(y))=\frac{1}{n}
	\;\;\Rightarrow\;\;
	\Pr(h(x)\neq h(y))=1-\frac{1}{n}
	\]

	\[
	\Pr(\text{no collision})
	=
	\prod_{i=1}^{m-1}\left(1-\frac{i}{n}\right)
	\]

	\[
	\ge
	\left(1-\frac{m}{n}\right)^m
	\]

	So
	\[
	\left(1-\frac{m}{n}\right)^m
	=
	\left(1-\frac{m}{c m^2}\right)^m
	=
	\left(1-\frac{1}{c m}\right)^m
	=
	\left(\left(1-\frac{1}{c m}\right)^{c m}\right)^{\frac{1}{c}}
	\]

	\[
	1-x \le e^{-x}
	\;\;\Leftrightarrow\;\;
	\frac{1}{2e}\le \left(1-\frac{1}{x}\right)^x \le \frac{1}{e},
	\qquad x=c m
	\]

	So
	\[
	\left(\left(1-\frac{1}{c m}\right)^{c m}\right)^{\frac{1}{c}}
	\ge
	\left(\frac{1}{2e}\right)^{\frac{1}{c}}
	\;\;\Rightarrow\;\;
	\Pr(\text{no collision}) \ge \left(\frac{1}{2e}\right)^{\frac{1}{c}}
	\]

	So
	\[
	\left(\frac{1}{2e}\right)^{\frac{1}{c}} \ge \frac{9}{10}
	\]

	\[
	\frac{1}{c}\ln\!\left(\frac{1}{2e}\right) \ge \ln\!\left(\frac{9}{10}\right)
	\]

	\[
	c \ge \frac{\ln(2e)}{\ln(10/9)} \approx 16.07
	\]

	If
	\[
	c > 16.07
	\quad\text{then}\quad
	\Pr(\text{no collision}) \ge \frac{9}{10}.
	\]

	\item (4 pts) A monkey types on a 26-letter keyboard that contains lowercase letters only. Each letter is chosen independently and uniformly at random from the alphabet. If the monkey types 100,000 letters, what is the expected number of times the sequence ``nyu'' appears? Give an upper bound on the probability that the sequence ``nyu'' appears at least ten times.

\[
\Pr(\text{NYU on 3 keys})=\left(\frac{1}{26}\right)^3
\;\Rightarrow\;
\Pr(\text{no NYU}) = 1-\left(\frac{1}{26}\right)^3
\]

\begin{itemize}
	\item So next question ``How many groups of 3 letters are there in $1\times 10^5$ random keys''. 
	\item Well each key is a ``start'' to a 3 letter pattern except last 2
	so $1\times 10^5 - 2$.
\end{itemize}

\[
\text{So } \mathbb{E}(\text{NYU appears in }1\times 10^5)
=
(1\times 10^5 - 2)\left(\frac{1}{26}\right)^3
\approx 5.6894
\]

Upper bound $\Rightarrow$ Markov inequality,
\[
\Pr(X>t)\le \frac{\mathbb{E}[X]}{t}
\]

\[
\Pr(\text{NYU appears} \ge 10)
\le
\frac{5.6894}{10}
=
0.56894
\]
\end{enumerate}

\section*{Problem 2: Why does Count-Min work so well in practice? (12 pts)}
We showed that Count-Min can estimate the frequency of any item in a stream of $n$ items up to additive error $\tfrac{m}{n}$ using $O(m)$ space. In practice it is often observed that this bound is pessimistic: the algorithm performs better than expected. In this problem, you will establish one reason why.

For any positive integer $m$, let $f_1,\dots,f_m$ be the frequencies of the $m$ most frequent items in our stream. Let
\[ C = n - \sum_{i=1}^m f_i. \]
In general, we can have that $C\ll n$. For example, it has been observed that up to 95\% of YouTube video views come from just 1\% of videos. Prove that using $O(m)$ space, Count-Min actually returns an estimate $\tilde f$ to $f(v)$ for any item $v$ satisfying
\[ f(v) \le \tilde f(v) \le f(v) + \frac{C}{m} \]
with $9/10$ probability. This is strictly better than the $\tfrac{m}{n}$ error bound shown in class.

\subsection*{Pr. 2}

\[
\text{In $U$, let } f_1,\ldots,f_m \text{ be the frequencies of the $m$ most frequent items.}
\]

\[
C = n - \sum_{i=1}^{m} f_i
= [\text{total stream in stream}] - [\text{sum of most frequent}]
\]

\[
C \ll n \;\;\Rightarrow\;\; \sum_{i=1}^{m} f_i = n - C \approx n
\]

Goal: prove in $O(m)$ Count-Min returns $\tilde f$ to $f(v)$ s.t.
\[
f(v) \le \tilde f \le f(v) + \frac{1}{m}C
\qquad \text{with prob. } \frac{9}{10},
\]
this is strictly better than the $\frac{1}{m}n$ error bound.

\bigskip

\textbf{Variables (count-min + my own)}
\begin{itemize}
	\item $n$: \# total items in stream
	\item $U$: Universe of all possible items
	\item $m$: \# of ``heavy hitters''
	\item $T$: the set of the tail i.e. not heavy hitters
	\item $C$: tail mass (sum of freq in $T$)
	\[ C=\sum_{x\in T} f(x) = n - \sum_{x\notin T} f(x) \]
	\item $w$: width of each Array[row] \quad (*think count-min from class: $w$ = \# indices, $d$ = \# arrays)
	\item $d$: \# of arrays [rows] \& hash functions
	\item[*] Each array has own hash function
\end{itemize}

\textbf{Configuration}\; \Big[ \text{This comes from a previous solve which I am simplifying} \Big]

To ensure use $O(m)$ space we set
\[ w=4m,\qquad d=4 \qquad\Longrightarrow\qquad 4m\cdot 4 = 16m = O(m) \]

\bigskip
\textbf{Error Analysis}

We state that the sketch is
\[ \tilde f(v)=\min_{j} A_j[h_j(v)]. \]
Since count-min doesn't under-estimate,
\[ \tilde f(v)\ge f(v). \]

Consider a row $j$:
\[ A_j[h_j(v)] = f(v) + \sum_{\substack{y\in U\\ y\neq v}} f(y)\cdot \mathbf{1}\{h_j(y)=h_j(v)\} \]

\[\text{Noise-Errors} =
\sum_{\substack{y\notin T\\ y\neq v}} f(y)\cdot \mathbf{1}\{h_j(y)=h_j(v)\}
+
\sum_{\substack{y\in T\\ y\neq v}} f(y)\cdot \mathbf{1}\{h_j(y)=h_j(v)\}
\]

A row $j$ is bad if high error rate
\[ \text{Error} > \frac{C}{m}. \]

This happens if either two events occur.

\textbf{\#1) High error in heavy hitters: if any heavy hitters collide \ldots}

\[ \Pr(\text{Any collision of ``heavyHitters''})\le \sum_{\substack{y\notin T\\ y\neq v}} \Pr(h_j(y)=h_j(v))
= \sum_{\substack{y\notin T\\ y\neq v}} \frac{1}{w} \le \frac{m}{w} = \frac{1}{4} \]

\textbf{\#2) High error rate in tail. Want it less than $\frac{C}{m}$.}

Let $E_T$ be the error rate in the tail.
\[ \mathbb{E}(E_T)=\frac{C}{w} \]

Using Markov's,
\[ \Pr\!\left(E_T \ge \frac{C}{m}\right) \le \frac{\mathbb{E}[E_T]}{C/m} = \frac{C/w}{C/m} = \frac{m}{w} = \frac{1}{4} \]

\textbf{Overall}

\[ \Pr(\text{High error Heavy hitters}) \le \frac{1}{4},
\qquad
\Pr(\text{High error tail noise}) \le \frac{1}{4} \]

So
\[ \Pr(\text{row is bad}) = \Pr(\text{High error Heavy hitters} \ \cup\ \text{High error tail noise}) \le \frac{1}{4}+\frac{1}{4} = \frac{1}{2} \]

So how many arrays (rows) are needed so that, with probability at least $9/10$, we have
\(\tilde f \le f(v)+\frac{C}{m}\)? Solve
\[
1-\left(\tfrac{1}{2}\right)^d = \tfrac{9}{10},
\]
which gives $d=4$. With $w=4m$ and $d=4$ we use $4\cdot 4m=16m=O(m)$ space. Hence the error bound holds with probability at least $9/10$. \(\Box\)

\bigskip
{\small *Possible risky notation spot: some sources use $n$ for universe size; here $n$ is total stream count.}


\section*{Problem 3: Randomized methods for efficient disease identification (group testing). (12 pts)}
One of the most important factors in controlling diseases like bird flu or, a few years ago, COVID-19, is testing. However, testing often requires processing in a lab, so can be expensive and slow. One way to make it cheaper is to test patients/livestock/etc.
in groups. The biological samples from multiple individuals (e.g., multiple nose swabs) are combined into a single test tube and tested for the disease all at once. If the test comes back negative, we know everyone in the group is negative. If the test comes back positive, we do not know which patients in the group actually have the disease, so further testing would be necessary. There's a trade-off here, but it turns out that, overall, group testing can save on the total number of tests run.

\begin{enumerate}[label=\arabic*.]
	\item (6 pts) Consider the following deterministic ``two-level'' testing scheme. We divide a population of $n$ individuals to be tested into $C$ groups of the same size. We then test each of these groups. For any group that comes back positive, we retest all members of the group individually. Show that there is a choice for $C$ such that, if $k$ individuals in the population have a disease (would test positive), we can find all of those individuals with $\le 2\sqrt{nk}$ tests. You can assume $k$ is known in advance (often it can be estimated accurately from the positive rate of prior tests). This is already an improvement on the naive $n$ tests when $k<0.25\,n$.

	\item (6 pts) We can use randomness to do better. Consider the following scheme: collect $q=O(\log n)$ biological samples from each individual (in reality, divide one sample into $q$ parts). Then, repeat the following process $q$ times: randomly partition our set of $n$ individuals into $C$ groups, and test each group in aggregate. Once this process is complete, report that an individual ``is positive'' if the groups they were part of tested positive all $q$ times. Report that an individual ``is negative'' if any of the groups they were part of tested negative. Prove that for $C=O(k)$, with probability $9/10$, this scheme finds all truly positive patients and reports no false positives. Thus, we only require $O(k\log n)$ tests!
\end{enumerate}

\section*{Problem 4: Try out mark-and-recapture! (12 pts)}
\begin{enumerate}[label=\arabic*.]
	\item (6 pts)

	% Pr. 4  (Mark-and-Recapture)  --- LaTeX for this section only

	% ---------------------------
	% Setup + Exploration
	% ---------------------------
	\[
	\text{Collecting } O\!\left(\frac{\sqrt{n}}{\epsilon}\right) \text{ samples (uniformly) from a set of size } n
	\text{, we can return } \tilde n \text{ such that, with } \Pr \ge \frac{9}{10},
	\]
	\[
	(1-\epsilon)n \le \tilde n \le (1+\epsilon)n.
	\]

	\[
	\mathbb{E}[D]
	= \sum_{1\le i<j\le m} \Pr(X_i=X_j)
	= \sum_{1\le i<j\le m} \frac{1}{n}
	= \binom{m}{2}\frac{1}{n}
	=: \mu,
	\]
	where \(m=\#\) queries (samples) and \(n=\#\) unique items. Also \(D=\#\) duplicates (collisions).

	We can flip this to define the estimator
	\[
	\tilde n := \frac{m(m-1)}{2D}.
	\]
	\[
	\text{(*) Note: this is a plug-in idea, replacing }\mathbb{E}[D]=\mu \text{ by the observed } D.
	\]

	% ---------------------------
	% Finding Var(D) \le \binom{m}{2}\frac{1}{n}
	% ---------------------------
	Define indicator variables, for \(1\le i<j\le m\),
	\[
	I_{ij} :=
	\begin{cases}
	1 & \text{if } X_i=X_j,\\
	0 & \text{otherwise}.
	\end{cases}
	\]
	Then
	\[
	D=\sum_{1\le i<j\le m} I_{ij},
	\qquad
	\mathbb{E}[I_{ij}] = \Pr(X_i=X_j)=\frac{1}{n}.
	\]

	We compute
	\[
	\operatorname{Var}(D)=\mathbb{E}[D^2]-\bigl(\mathbb{E}[D]\bigr)^2.
	\]
	Also,
	\[
	D^2
	=
	\left(\sum_{i<j} I_{ij}\right)^2
	=
	\sum_{i<j} I_{ij}^2
	+
	2\sum_{\substack{(i<j),(k<\ell)\\(i<j)<(k<\ell)}} I_{ij}I_{k\ell}.
	\]
	Taking expectations,
	\[
	\mathbb{E}[D^2]
	=
	\sum_{i<j}\mathbb{E}[I_{ij}^2]
	+
	2\sum_{(i<j)<(k<\ell)} \mathbb{E}[I_{ij}I_{k\ell}].
	\]
	Since \(I_{ij}^2=I_{ij}\) (Bernoulli),
	\[
	\mathbb{E}[I_{ij}^2]=\mathbb{E}[I_{ij}]=\frac{1}{n}.
	\]

	Now write variance in the standard way:
	\[
	\operatorname{Var}(D)
	=
	\sum_{i<j}\operatorname{Var}(I_{ij})
	+
	2\sum_{(i<j)<(k<\ell)} \operatorname{Cov}(I_{ij},I_{k\ell}).
	\]
	We treat the covariance terms as \(0\) (as in the notes we used), hence
	\[
	\operatorname{Var}(D)=\sum_{i<j}\operatorname{Var}(I_{ij}).
	\]
	Each \(I_{ij}\) is Bernoulli with \(p=\frac{1}{n}\), so
	\[
	\operatorname{Var}(I_{ij}) = p(1-p)\le p=\frac{1}{n}.
	\]
	Therefore,
	\[
	\operatorname{Var}(D)
	\le \sum_{i<j}\frac{1}{n}
	=\binom{m}{2}\frac{1}{n}
	=\mu.
	\]

	% ---------------------------
	% Chebyshev + Translating to a bound on \tilde n
	% ---------------------------
	Chebyshev (variance form):
	\[
	\Pr\bigl(|D-\mu|\ge t\bigr)\le \frac{\operatorname{Var}(D)}{t^2}.
	\]
	Using \(\operatorname{Var}(D)\le \mu\),
	\[
	\Pr\bigl(|D-\mu|\ge t\bigr)\le \frac{\mu}{t^2}.
	\]
	So by complement,
	\[
	\Pr\bigl(|D-\mu|\le t\bigr)\ge 1-\frac{\mu}{t^2}.
	\]

	Set \(t=\alpha\mu\) (with \(\alpha>0\)):
	\[
	\Pr\bigl(|D-\mu|\le \alpha\mu\bigr)\ge 1-\frac{1}{\alpha^2\mu}.
	\]
	The event \(|D-\mu|\le \alpha\mu\) is equivalent to
	\[
	(1-\alpha)\mu \le D \le (1+\alpha)\mu.
	\]
	Substitute \(\mu=\binom{m}{2}\frac{1}{n}\):
	\[
	(1-\alpha)\binom{m}{2}\frac{1}{n} \le D \le (1+\alpha)\binom{m}{2}\frac{1}{n}.
	\]
	Invert (all terms are positive) and use \(\tilde n=\frac{\binom{m}{2}}{D}\):
	\[
	\frac{n}{1+\alpha}\le \tilde n \le \frac{n}{1-\alpha}.
	\]

	% ---------------------------
	% Choose parameters for probability 9/10, then connect to \epsilon
	% ---------------------------
	To get probability at least \(9/10\), it suffices that
	\[
	1-\frac{1}{\alpha^2\mu}\ge \frac{9}{10}
	\quad\Longleftrightarrow\quad
	\frac{1}{\alpha^2\mu}\le \frac{1}{10}
	\quad\Longleftrightarrow\quad
	\mu \ge \frac{10}{\alpha^2}.
	\]
	Since \(\mu=\binom{m}{2}\frac{1}{n}\),
	\[
	\binom{m}{2}\frac{1}{n}\ge \frac{10}{\alpha^2}
	\quad\Longleftrightarrow\quad
	\binom{m}{2}\ge \frac{10n}{\alpha^2}.
	\]

	To match the target \((1-\epsilon)n\le \tilde n\le (1+\epsilon)n\), it is enough to choose
	\[
	\alpha \le \frac{\epsilon}{1+\epsilon}.
	\]
	A simple safe choice is \(\alpha=\epsilon/2\), which yields
	\[
	m = O\!\left(\frac{\sqrt{n}}{\epsilon}\right)
	\quad \text{and} \quad
	\Pr\!\left((1-\epsilon)n\le \tilde n \le (1+\epsilon)n\right)\ge \frac{9}{10}.
	\]
\end{enumerate}

\section*{Problem 5: Distributed Importance Sampling. (15 pts)}
In many machine learning training pipelines, not all data points are created equal. We often want to perform Weighted Reservoir Sampling: given a stream of items $x_1,x_2,\dots$ with associated positive weights $w_1,w_2,\dots$, we want to maintain a sample of size $k$. The goal is that, at any point in time, the probability that item $x_i$ is included in the sample is proportional to its weight relative to the other items seen so far.

A naive approach is difficult because standard reservoir sampling relies on uniform probabilities. However, in 2006, Efraimidis and Spirakis proposed a specialized algorithm suitable for distributed systems. The Algorithm:
\begin{itemize}
	\item For each item $x_i$ arriving in the stream with weight $w_i$, generate a random number $u_i$ uniformly from $(0,1]$.
	\item Compute a ``key'' for the item: $k_i = u_i^{1/w_i}$.
	\item Maintain the $k$ items with the largest keys $k_i$ seen so far.
\end{itemize}

\begin{enumerate}[label=(\alph*)]
	\item (5 pts) The Single Item Case ($k=1$). Consider just two items $x_1,x_2$ with weights $w_1,w_2$. Prove that the probability that $x_1$ has the larger key (and is thus selected) is exactly $\dfrac{w_1}{w_1+w_2}$. Hint: you are comparing two random variables $K_1=U_1^{1/w_1}$ and $K_2=U_2^{1/w_2}$. Set up the integral $\int_0^1 \Pr[U_1^{1/w_1}>y] f_{K_2}(y)\,dy$ or a similar double integral.

	\item (5 pts) Generalizing to Stream. Using your result from part (a), argue why this method works for a stream of $n$ items for $k=1$. Specifically, show that for any item $i$, the probability it has the maximum key is $\dfrac{w_i}{\sum_{j=1}^n w_j}$.

	\item (5 pts) Distributed Merging. One of the massive advantages of this specific algorithm is that it is ``mergeable.'' Suppose we have two servers, A and B. Server A observes a stream of $n_A$ items and maintains a weighted reservoir sample $S_A$ of size $k$ using the algorithm above. Server B observes a distinct stream of $n_B$ items and maintains $S_B$, also of size $k$. Describe a procedure to merge $S_A$ and $S_B$ into a single global reservoir sample $S_{\text{global}}$ of size $k$ that represents the weighted sample of the union of both streams. Prove that your merged sample has the same distribution as if a single server had processed all $n_A+n_B$ items sequentially.
\end{enumerate}

\end{document}
